\chapter{Local Cohomology and Duality}

In this chatper, assume all rings are Noetherian.

\section{Local cohomology}

Let \(R\) be a commutative ring, \(I\subset R\) is an ideal and \(M\) is an \(R\)-module. 

\begin{definition}[Local Cohomology]
     We define a submodule of \(M\) as 
     \[\Gamma_I(M)=\left\{ m\in M\mid I^nm=0\ \ for\ some\ n\in \mathbb{N} \right\}\]
     When \(\Gamma_I(M)=M\), we say \(M\) is said to be \(I\)-torsion. For any homomorphism of \(R\)-modules \(\varphi:M\rightarrow N\), we have an induced homomorphism
     \[\Gamma_I(\varphi):\Gamma_I(M)\rightarrow \Gamma_I(N)\]
     So \(\Gamma_I(-)\) is a functor on the category of \(R\)-modules, and it can be extended to complexes of \(R\)-modules.

     It is not hard to prove that \(\Gamma_I(-)\) is left exact, we define 
     \[H^k_I(-):=R^k\Gamma_I(-)\]
     the \(k\)th right derived functor of \(\Gamma_I(-)\), which is called \(k\)th local cohomology with support in \(I\).
\end{definition}

Here are a few basic properties of local cohomology:

\begin{proposition} 
     Let \(R\) be a ring, \(I\subset R\) an ideal and \(M\) an \(R\)-module.
     \begin{enumerate}[(1)]
       \item One has \(H_I^0(M)\cong \Gamma_I(M)\), and \(H^k_I(M)\) is \(I\)-torsion for all \(k\).
       \item If \(\sqrt{I}=\sqrt{J}\), then \(H_I^k(M)\cong H_J^k(M)\) for all \(k\).
       \item Let \(\left\{ M_\lambda \right\}\) be a family of \(R\)-modules. For each integer \(k\), we have 
       \[H^k_I(\bigoplus_\lambda M_\lambda)\cong \bigoplus_\lambda H^k_I(M_\lambda).\]
       \item Given a short exact sequence of \(R\)-modules 
       \[0\rightarrow M'\rightarrow M\rightarrow M''\rightarrow 0,\]
       we have a long exact sequence in local cohomology
       \[\cdots\rightarrow H^k_I(M')\rightarrow H^k_I(M)\rightarrow H^k_I(M'')\rightarrow H^{k+1}_I(M')\rightarrow \cdots\]
     \end{enumerate}
\end{proposition}

\begin{proof}
    \begin{enumerate}[(1)]
      \item By taking an injective rresolution.
      \item Note that \(I\subset \sqrt{I}=\sqrt{J}\) and \(I\) is finitely generated (because \(R\) is Noetherian), we have \(I^n\subset J\) for some \(n\). 
      \item \(\Gamma_I(-)\) and taking homology commute with direct sum.
      \item Standard homological algebra.
    \end{enumerate}
\end{proof}

\begin{example}
    Let \(p\) be a prime number and \(I=(p)\subset \mathbb{Z}\). We compute \(H_I^j(\mathbb{Z})\). Consider the injective resolution 
    \[0\rightarrow \mathbb{Z}\hookrightarrow \mathbb{Q}\rightarrow \mathbb{Q}/\mathbb{Z}\rightarrow 0.\]
    All \(H^j_I(\mathbb{Z})\) vanishes except for \(j=1\). Note that \(\Gamma_I(\mathbb{Q})=0\) and \(\Gamma_I(\mathbb{Q}/\mathbb{Z})=\mathbb{Z}[p^{-1}]/\mathbb{Z}\). So \(H^1_I(\mathbb{Z})=\mathbb{Z}[p^{-1}]/\mathbb{Z}\).
\end{example}

\section{Other ways to compute local cohomology}

There are other ways to compute local cohomology.

\begin{theorem}
  For each \(R\)-module \(M\), there are natural isomorphisms
  \[\varinjlim_n \ext^j_R(R/I^n,M)\cong H^j_I(M)\]
  for each \(j\geq 0\).
\end{theorem}

\begin{proof}
    For each \(R\)-module \(E\) and each positive integer \(n\), we have a functorial isomorphism
    \begin{align*}
         \hom_R(R/I^n,E)&\xrightarrow{\cong} \left\{ x\in E\mid I^n\cdot x=0 \right\}\\
         f&\mapsto f(1).
    \end{align*}
\end{proof}

\begin{remark}
     Let \(\left\{ I_n \right\}_{n\geq 0}\) be a decreasing chain of ideals cofinal with the chain of ideals \(\left\{ I^n \right\}_{n\geq 0}\). That is, for any \(n\geq 0\), there exists positive integers \(p,q>0\) such that \(I^{n+p}\subseteq I_n \) and \(I_{n+q}\subseteq I^n\). Recall that both can be used to compute the directed limit, so we get a functorial isomorphism 
     \[\varinjlim_n \ext^j_R(R/I_n,M)\cong H^j_I(M).\]
\end{remark}

We can also use the Koszul complex to compute the local cohomology. Recall that for an element \(x\in R\) and \(n\in \mathbb{N}\), the Koszul complex \(K^\bullet(x^n;R)\) is 
\[0\rightarrow R\xrightarrow{x^n} R\rightarrow 0\]
where the nonzero term is in degree -1 and 0. We have a map of Koszul complexes \(K^\bullet(x^{n+1};R)\rightarrow K^\bullet(x^n;R)\) given by 
% https://q.uiver.app/#q=WzAsOCxbMCwwLCIwIl0sWzEsMCwiUiJdLFsyLDAsIlIiXSxbMywwLCIwIl0sWzAsMSwiMCJdLFsxLDEsIlIiXSxbMiwxLCJSIl0sWzMsMSwiMCJdLFswLDFdLFsxLDIsInhee24rMX0iXSxbMiwzXSxbNSw2LCJ4Xm4iXSxbNCw1XSxbNiw3XSxbMSw1LCJ4IiwyXSxbMiw2LCJpZCJdXQ==
\[\begin{tikzcd}
    0 & R & R & 0 \\
    0 & R & R & 0
    \arrow[from=1-1, to=1-2]
    \arrow["{x^{n+1}}", from=1-2, to=1-3]
    \arrow["x"', from=1-2, to=2-2]
    \arrow[from=1-3, to=1-4]
    \arrow["id", from=1-3, to=2-3]
    \arrow[from=2-1, to=2-2]
    \arrow["{x^n}", from=2-2, to=2-3]
    \arrow[from=2-3, to=2-4]
  \end{tikzcd}\]
It is not hard to see that \(\left\{ K^\bullet(x^n;R) \right\}_{n\geq 1}\) form an inverse system. More generally, set \(\underline{x}^n\) is a sequence \(x_1^n,\ldots,x_d^n\). The Koszul complex \(K^\bullet(\underline{x}^n;R)\) is given by 
\[K^\bullet(\underline{x}^n;R)=K^\bullet(x_1^n;R)\otimes_R\cdots\otimes_RK^\bullet(x_d^n;R).\]
We obtain an inverse system of complexes of \(R\)-modules 
\[\cdots\rightarrow K^\bullet(\underline{x}^{n+1};R)\rightarrow K^\bullet(\underline{x}^n;R)\rightarrow \cdots\]
Note that for every \(n\), we have a morphism
\[K^\bullet(\underline{x}^n;R)\rightarrow R/(x_1^n,\ldots,x_d^n).\]

\section{Depth and cohomological dimension}
